% Options for packages loaded elsewhere
% Options for packages loaded elsewhere
\PassOptionsToPackage{unicode}{hyperref}
\PassOptionsToPackage{hyphens}{url}
%
\documentclass[
  11pt,
  letterpaper,
  DIV=11,
  numbers=noendperiod,
  twoside]{scrartcl}
\usepackage{xcolor}
\usepackage[left=1.8in, right=1.8in, top=1.5in, bottom=1.5in,
includemp=TRUE, marginparwidth=0in, marginparsep=0in]{geometry}
\usepackage{amsmath,amssymb}
\setcounter{secnumdepth}{3}
\usepackage{iftex}
\ifPDFTeX
  \usepackage[T1]{fontenc}
  \usepackage[utf8]{inputenc}
  \usepackage{textcomp} % provide euro and other symbols
\else % if luatex or xetex
  \usepackage{unicode-math} % this also loads fontspec
  \defaultfontfeatures{Scale=MatchLowercase}
  \defaultfontfeatures[\rmfamily]{Ligatures=TeX,Scale=1}
\fi
\usepackage{lmodern}
\ifPDFTeX\else
  % xetex/luatex font selection
  \setmathfont[]{Garamond-Math}
\fi
% Use upquote if available, for straight quotes in verbatim environments
\IfFileExists{upquote.sty}{\usepackage{upquote}}{}
\IfFileExists{microtype.sty}{% use microtype if available
  \usepackage[]{microtype}
  \UseMicrotypeSet[protrusion]{basicmath} % disable protrusion for tt fonts
}{}
\usepackage{setspace}
% Make \paragraph and \subparagraph free-standing
\makeatletter
\ifx\paragraph\undefined\else
  \let\oldparagraph\paragraph
  \renewcommand{\paragraph}{
    \@ifstar
      \xxxParagraphStar
      \xxxParagraphNoStar
  }
  \newcommand{\xxxParagraphStar}[1]{\oldparagraph*{#1}\mbox{}}
  \newcommand{\xxxParagraphNoStar}[1]{\oldparagraph{#1}\mbox{}}
\fi
\ifx\subparagraph\undefined\else
  \let\oldsubparagraph\subparagraph
  \renewcommand{\subparagraph}{
    \@ifstar
      \xxxSubParagraphStar
      \xxxSubParagraphNoStar
  }
  \newcommand{\xxxSubParagraphStar}[1]{\oldsubparagraph*{#1}\mbox{}}
  \newcommand{\xxxSubParagraphNoStar}[1]{\oldsubparagraph{#1}\mbox{}}
\fi
\makeatother


\usepackage{longtable,booktabs,array}
\newcounter{none} % for unnumbered tables
\usepackage{calc} % for calculating minipage widths
% Correct order of tables after \paragraph or \subparagraph
\usepackage{etoolbox}
\makeatletter
\patchcmd\longtable{\par}{\if@noskipsec\mbox{}\fi\par}{}{}
\makeatother
% Allow footnotes in longtable head/foot
\IfFileExists{footnotehyper.sty}{\usepackage{footnotehyper}}{\usepackage{footnote}}
\makesavenoteenv{longtable}
\usepackage{graphicx}
\makeatletter
\newsavebox\pandoc@box
\newcommand*\pandocbounded[1]{% scales image to fit in text height/width
  \sbox\pandoc@box{#1}%
  \Gscale@div\@tempa{\textheight}{\dimexpr\ht\pandoc@box+\dp\pandoc@box\relax}%
  \Gscale@div\@tempb{\linewidth}{\wd\pandoc@box}%
  \ifdim\@tempb\p@<\@tempa\p@\let\@tempa\@tempb\fi% select the smaller of both
  \ifdim\@tempa\p@<\p@\scalebox{\@tempa}{\usebox\pandoc@box}%
  \else\usebox{\pandoc@box}%
  \fi%
}
% Set default figure placement to htbp
\def\fps@figure{htbp}
\makeatother


% definitions for citeproc citations
\NewDocumentCommand\citeproctext{}{}
\NewDocumentCommand\citeproc{mm}{%
  \begingroup\def\citeproctext{#2}\cite{#1}\endgroup}
\makeatletter
 % allow citations to break across lines
 \let\@cite@ofmt\@firstofone
 % avoid brackets around text for \cite:
 \def\@biblabel#1{}
 \def\@cite#1#2{{#1\if@tempswa , #2\fi}}
\makeatother
\newlength{\cslhangindent}
\setlength{\cslhangindent}{1.5em}
\newlength{\csllabelwidth}
\setlength{\csllabelwidth}{3em}
\newenvironment{CSLReferences}[2] % #1 hanging-indent, #2 entry-spacing
 {\begin{list}{}{%
  \setlength{\itemindent}{0pt}
  \setlength{\leftmargin}{0pt}
  \setlength{\parsep}{0pt}
  % turn on hanging indent if param 1 is 1
  \ifodd #1
   \setlength{\leftmargin}{\cslhangindent}
   \setlength{\itemindent}{-1\cslhangindent}
  \fi
  % set entry spacing
  \setlength{\itemsep}{#2\baselineskip}}}
 {\end{list}}
\usepackage{calc}
\newcommand{\CSLBlock}[1]{\hfill\break\parbox[t]{\linewidth}{\strut\ignorespaces#1\strut}}
\newcommand{\CSLLeftMargin}[1]{\parbox[t]{\csllabelwidth}{\strut#1\strut}}
\newcommand{\CSLRightInline}[1]{\parbox[t]{\linewidth - \csllabelwidth}{\strut#1\strut}}
\newcommand{\CSLIndent}[1]{\hspace{\cslhangindent}#1}



\setlength{\emergencystretch}{3em} % prevent overfull lines

\providecommand{\tightlist}{%
  \setlength{\itemsep}{0pt}\setlength{\parskip}{0pt}}



 


\setlength\heavyrulewidth{0ex}
\setlength\lightrulewidth{0ex}
\usepackage[automark]{scrlayer-scrpage}
\clearpairofpagestyles
\cehead{
  Brian Weatherson
  }
\cohead{
  New Work for Dempster-Shafer Functions
  }
\ohead{\bfseries \pagemark}
\cfoot{}
\makeatletter
\newcommand*\NoIndentAfterEnv[1]{%
  \AfterEndEnvironment{#1}{\par\@afterindentfalse\@afterheading}}
\makeatother
\NoIndentAfterEnv{itemize}
\NoIndentAfterEnv{enumerate}
\NoIndentAfterEnv{description}
\NoIndentAfterEnv{quote}
\NoIndentAfterEnv{equation}
\NoIndentAfterEnv{longtable}
\NoIndentAfterEnv{abstract}
\renewenvironment{abstract}
 {\vspace{-1.25cm}
 \quotation\small\noindent\emph{Abstract}:}
 {\endquotation}
\newfontfamily\tfont{EB Garamond}
\addtokomafont{disposition}{\rmfamily}
\addtokomafont{descriptionlabel}{\rmfamily}
\addtokomafont{title}{\normalfont\itshape}
\let\footnoterule\relax

\makeatletter
\renewcommand{\@maketitle}{%
  \newpage
  \null
  \vskip 2em%
  \begin{center}%
  \let \footnote \thanks
    {\itshape\huge\@title \par}%
    \vskip 0.5em%  % Reduced from default
    {\large
      \lineskip 0.3em%  % Reduced from default 0.5em
      \begin{tabular}[t]{c}%
        \@author
      \end{tabular}\par}%
    \vskip 0.5em%  % Reduced from default
    {\large \@date}%
  \end{center}%
  \par
  }
\makeatother
\RequirePackage{lettrine}

\renewenvironment{abstract}
 {\quotation\small\noindent\emph{Abstract}:}
 {\endquotation\vspace{-0.02cm}}

\setmainfont{EB Garamond Math}[
  BoldFont = {EB Garamond SemiBold},
  ItalicFont = {EB Garamond Italic},
  RawFeature = {+smcp},
]

\newfontfamily\scfont{EB Garamond Regular}[RawFeature=+smcp]
\renewcommand{\textsc}[1]{{\scfont #1}}

\renewcommand{\LettrineTextFont}{\scfont}
\newcommand{\nmodels}{\mathrel{\ooalign{$\models$\cr\raisebox{-0.001ex}{\hss$\mkern-1mu/\hss$}}}}
\newcommand{\llbracket}{[\![}
\newcommand{\rrbracket}{]\!]}
\usepackage{pifont}
\newcommand{\tick}{\text{\ding{51}}}
\newcommand{\cross}{\text{\ding{55}}}
\setlength\heavyrulewidth{0ex}
\setlength\lightrulewidth{0ex}
\usepackage[lines=2]{lettrine}
\cehead{Draft of February 28, 2026}
\DeclareSymbolFont{symbolsC}{U}{txsyc}{m}{n}
\DeclareMathSymbol{\Boxright}{\mathrel}{symbolsC}{128}
\cehead{
       Brian Weatherson
       }
\KOMAoption{captions}{tableheading}
\makeatletter
\@ifpackageloaded{caption}{}{\usepackage{caption}}
\AtBeginDocument{%
\ifdefined\contentsname
  \renewcommand*\contentsname{Table of contents}
\else
  \newcommand\contentsname{Table of contents}
\fi
\ifdefined\listfigurename
  \renewcommand*\listfigurename{List of Figures}
\else
  \newcommand\listfigurename{List of Figures}
\fi
\ifdefined\listtablename
  \renewcommand*\listtablename{List of Tables}
\else
  \newcommand\listtablename{List of Tables}
\fi
\ifdefined\figurename
  \renewcommand*\figurename{Figure}
\else
  \newcommand\figurename{Figure}
\fi
\ifdefined\tablename
  \renewcommand*\tablename{Table}
\else
  \newcommand\tablename{Table}
\fi
}
\@ifpackageloaded{float}{}{\usepackage{float}}
\floatstyle{ruled}
\@ifundefined{c@chapter}{\newfloat{codelisting}{h}{lop}}{\newfloat{codelisting}{h}{lop}[chapter]}
\floatname{codelisting}{Listing}
\newcommand*\listoflistings{\listof{codelisting}{List of Listings}}
\makeatother
\makeatletter
\makeatother
\makeatletter
\@ifpackageloaded{caption}{}{\usepackage{caption}}
\@ifpackageloaded{subcaption}{}{\usepackage{subcaption}}
\makeatother
\usepackage{bookmark}
\IfFileExists{xurl.sty}{\usepackage{xurl}}{} % add URL line breaks if available
\urlstyle{same}
% fallback for those not using the hyperref driver hyperxmp:
\makeatletter
\@ifundefined{xmpquote}{\newcommand{\xmpquote}[1]{#1}}{}
\makeatother
\hypersetup{
  pdftitle={New Work for Dempster-Shafer Functions},
  pdfauthor={Brian Weatherson},
  hidelinks,
  pdfcreator={LaTeX via pandoc}}


\title{New Work for Dempster-Shafer Functions}
\author{Brian Weatherson}
\date{2026}
\begin{document}
\maketitle
\begin{abstract}
Works of fiction typically raise questions they do not answer. I argue
that some of these questions have no answers. There is no fact of the
matter about just when some character was born, or what their heart rate
was at any given moment. This seems intuitive, almost obvious, but there
are arguments against it. Kyle Blumberg and Ben Holguín have recently
argued that the way Lewis captures this intuition leads to problems. I
show that we avoid the problems if we model fictions using Humberstone's
possibility frames. That modelling choice has a striking consequence for
the theory of credence. Uncertainty about fictions, and uncertainty
about any topic where there might not be a fact of the matter, need not
be probabilistic. The accuracy-dominance argument for probabilism
presupposes bivalence; it requires that every question has an answer.
The considerations behind the accuracy-dominance argument are still
sound; it's just that without bivalence, they don't lead to probabilism.
Rather, they imply that credences should be modelled by Dempster-Shafer
belief functions. This is so even if the logic of the fiction, and of
our talk about the fiction, is resolutely classical.
\end{abstract}


\setstretch{1.1}
Eleanor Rigby, we are told, picked up the rice in a church where a
wedding had been. Have you ever thought about that wedding? Who were the
happy couple? What were their names? Did they stay together for years,
eventually bounce grandchildren on their knees? The song says enough to
make it possible to raise these questions, but not to answer them.

\emph{Eleanor Rigby}, the song, is special in many respects, but this is
not one of them. Works of fiction typically raise questions they do not
answer. Mostly these questions are not worth even trying to answer,
though in some cases at least trying to answer them feels important to
engaging with the fiction. (\emph{Blade Runner} is a famous example of
this.) Still, even the silly questions can be philosophically
interesting. For while it's not important what the answers to those
questions are, it is somewhat important what an answer to the questions
looks like. That's because thinking about questions about fictions that
may not have answers serves as an exemplar for modelling uncertainty
about any realm where there might not be a fact of the matter. That's
our ultimate subject, and we'll get there by looking in some detail at
the fictional case.

What's true in a fiction is sometimes settled by the words of the song.
In realistic fictions, others are settled by what would have to be true
for the text to make sense. For instance, the married couple in
\emph{Eleanor Rigby} (hereafter, The Couple) must have names. How else
would Father McKenzie marry them? Some are settled by the context the
fiction was produced in. \emph{Eleanor Rigby} was released in 1966, and
it feels like it's set in the reasonably recent past, so the wedding
takes place probably in the 1960s, or perhaps the late 1950s. It's not
meant to take place too far away, so it's certainly set in the UK, and
probably in England. Others are settled by facts about the world which
are imported in. For instance, we know Rocky Raccoon lives in
\emph{South} Dakota (at least while he lives) because he lives in the
Black Hills, and these are actually in South Dakota. Relatedly, because
all church marriages in the UK in the 1960s were heterosexual, The
Couple must consist of The Husband and The Wife. Perhaps the intent of
the author also affects what's true in the story, beyond the factors
already mentioned; we'll leave that question to one side.

What we will focus on is whether these questions must have answers. Is
there a fact about what The Husband's name is? A fairly orthodox theory,
developed in most detail by David Lewis
(\citeproc{ref-Lewis1978b}{1978}), says that there is not. Lewis's
theory has three parts. First, sentences like (1) and (2), have an
implicit operator in them.

\begin{enumerate}
\def\labelenumi{(\arabic{enumi})}
\tightlist
\item
  The Husband's name is Vincent.
\item
  The Husband has a name.
\end{enumerate}

We should understand these as saying

\begin{enumerate}
\def\labelenumi{(\arabic{enumi})}
\setcounter{enumi}{2}
\tightlist
\item
  In the fiction, The Husband's name is Vincent.
\item
  In the fiction, The Husband has a name.
\end{enumerate}

Second, the logic of this operator is that of a box in a normal modal
logic which is not \emph{functional} in the sense described by Shawn
Standefer (\citeproc{ref-Standefer2018}{2018}). In particular, it does
not include (5) as an axiom.\footnote{Standefer is interested primarily
  in the logic which also has \(\Diamond \top\) as an axiom. The
  designation of these logics as `functional' makes sense in a possible
  worlds framework, but is somewhat misleading in the possibility frame
  setting of this paper.}

\begin{enumerate}
\def\labelenumi{(\arabic{enumi})}
\setcounter{enumi}{4}
\tightlist
\item
  \(\Box \Phi \vee \Box \neg \Phi\)
\end{enumerate}

For Lewis, there are many worlds which have an equally good claim to be
the world of the fiction, and something is only true in the fiction if
it is true at all of them. So the English sentence (6) is ambiguous
between (7) and (8), with (7) false and (8) true. Because (3) and (7)
are both false, (5) fails.

\begin{enumerate}
\def\labelenumi{(\arabic{enumi})}
\setcounter{enumi}{5}
\tightlist
\item
  The Husband's name is not Vincent.
\item
  In the fiction, The Husband's name is not Vincent.
\item
  It is not the case that in the fiction, The Husband's name is Vincent.
\end{enumerate}

Third, Lewis thinks that this `in the fiction' operator can be
understood in terms of counterfactuals. Roughly, what's true in the
fiction is what would be true were the text of the story told as known
truth.

In this paper, I'll accept the first of Lewis's claims (that there is an
operator), reject the second (that it has a non-functional logic), and
mostly set aside the third. In `Fictional Reality', Kyle Blumberg and
Ben Holguín (\citeproc{ref-BlumbergHolguin2025}{2025}) run a series of
arguments against this second assumption of Lewis's. I'll briefly
mention a handful that I find most compelling.

For many questions that seem to be left open by a fiction, it makes
sense to wonder about them. This can range from questions that are
central to how one understands a fiction, like \emph{Is Deckard a
replicant?}, to trivial questions like \emph{Do The Couple have
grandchildren?}. But on Lewis's picture, there isn't anything to wonder
about. We know that it's not true that, in the fiction, they have
grandchildren, and it's also not true that, in the fiction, they do not
have grandchildren. Intuitively, (9) is true, but there isn't any
sensible way to insert an \emph{In the fiction} operator that makes it
true on Lewis's theory.

\begin{enumerate}
\def\labelenumi{(\arabic{enumi})}
\setcounter{enumi}{8}
\tightlist
\item
  I don't know whether The Couple had grandchildren.
\end{enumerate}

It isn't defective to wonder about this, in the way it is defective to
wonder about things that are settled by the fiction.

\begin{enumerate}
\def\labelenumi{(\arabic{enumi})}
\setcounter{enumi}{9}
\tightlist
\item
  I wonder whether The Couple had grandchildren.
\item
  \#I wonder whether Sherlock Holmes met Napoleon.
\end{enumerate}

For Lewis, once we know all there is to know about the fiction, all
worlds are sorted into two buckets; those that are consistent with the
fiction and those that are not. There is no way to make sense of
relative likelihoods within those buckets. So we can't make sense of
(12), but it feels intuitively true.

\begin{enumerate}
\def\labelenumi{(\arabic{enumi})}
\setcounter{enumi}{11}
\tightlist
\item
  The Couple might have had five grandchildren, and they might have had
  ten, but given the social norms at the time, and in the generation
  following, it's more likely they had five.
\end{enumerate}

Blumberg and Holguín infer from facts like these that there is a fact of
the matter about any question raised by the fiction. As they note, this
strikes many people as absurd. I'm one of those people, so I'd like to
set out a theory in which there need not be any such facts, without
falling into the traps besetting Lewis's theory.

One might wonder whether we need such detail on trivial questions like
the future of The Couple. This concern is in part a function of my
choosing a trivial question; many more natural questions about fictions
have the same shape. Who hasn't wondered what becomes of the
broken-hearted? More generally, the logical and epistemological issues
that this topic raises have implications far beyond fiction. Any time we
can coherently wonder about something without being certain whether
there is a fact of the matter, the same issues will arise. Our subject
is nothing less than the epistemology of the not clearly factual.

We can have graded uncertainty about questions where we aren't sure
there is a fact of the matter. This raises questions about the rational
constraints on this uncertainty. There are several familiar arguments
that all uncertainty should, ultimately, be probabilistic. I'll argue
that these arguments rest on an assumption that all the questions we're
uncertain about are (certainly) about matters of fact. When there might
be no fact of the matter, it is rational for one's uncertainty to be
represented not by a probability function, but by a Dempster-Shafer
belief function. (This is the `New Work' for these functions the title
promised.)

This paper has a snake-like structure. In sections \ref{sec-sec1} and
\ref{sec-sec5} I discuss \emph{possibility frames}. This is a system,
introduced by Lloyd Humberstone (\citeproc{ref-Humberstone1981a}{1981})
and recently discussed in detail by Wesley Holliday
(\citeproc{ref-Holliday2025}{2025}), which lets us preserve something
like classical logic while allowing that there might not be a fact of
the matter about some questions. I'll introduce the theory in section
\ref{sec-sec1}, and return to the slightly vexed question of whether
this is really \emph{classical} logic in section \ref{sec-sec5}. In
sections \ref{sec-sec2} and \ref{sec-sec4} I'll discuss Dempster-Shafer
functions. In section \ref{sec-sec2} I'll set out what they are, and the
standard arguments against using them. In section \ref{sec-sec4} I'll
come back to why these arguments aren't compelling when the subject
matter might be non-factual. In the middle, I'll come back to Blumberg
and Holguín's arguments about fiction, and defend the position already
sketched, that while their objections to Lewis are successful, their
proposed remedy is implausible and should be rejected.

\section{Possibility Frames}\label{sec-sec1}

A Humberstone possibility frame (hereafter, just \emph{possibility
frame}) is a quadruple \(\langle W, {\leqslant}, \mathbb{R}, V\rangle\),
where \(W\) is a set of possibilities, \(\leqslant\) is a
\emph{refinement} relation on them, \(\mathbb{R}\) is a countable family
of accessibility relations (with \(R_i\) as the \(i\)'th member), and
\(V\) is a valuation function.\footnote{One might expand frames even
  further, e.g., by adding a selection function that enters into the
  semantics for counterfactuals. Counterfactuals will be sitting just
  under the surface at many points in this paper, but to keep things
  focussed I'll defer discussion of counterfactuals to another day.}

One way to think about the possibilities in \(W\), suggested by
Humberstone (\citeproc{ref-Humberstone1981a}{1981}) and used extensively
by Holliday (\citeproc{ref-Holliday2025}{2025}), is that they are sets
of possible worlds. I prefer to think of them as \emph{stories}, at
least on the assumption that stories are incomplete. A possibility, like
a story, is something which settles some facts, while possibly leaving
others open.

Following Humberstone, I'll write \(x \leqslant y\) to mean that \(y\)
is a \textbf{refinement} of \(x\).\footnote{\(y \geqslant x\) will just
  mean the same thing as \(x \leqslant y\).} To say \(y\) refines \(x\)
is to say that \(y\) makes true everything that \(x\) makes true, and
possibly some more things besides. If possibilities are sets of possible
worlds, then \(y\) is a subset of \(x\). (This makes the notation here
confusing, which is why Holliday (\citeproc{ref-Holliday2025}{2025})
prefers to write \(y \sqsubseteq x\) when \(y\) is a refinement of
\(x\).) If possibilities are stories, then \(y\) is a refinement of
\(x\) if \(y\) fills in part of what \(x\) leaves open. This is how
I'll, intuitively, think of refinements.

\(V\) is a \textbf{valuation} function. It tells us where the atomics
are true. Formally, it is a function from \(\mathcal{P}\), the set of
atomics, to \(\wp(W)\), satisfying these two constraints:

\begin{itemize}
\tightlist
\item
  For all \(x\), if \(x \in V(p)\) and \(x \leqslant y\), then
  \(y \in V(p)\). Intuitively, truth for atomics is \textbf{persistent}
  across refinements.
\item
  For all \(x\), if
  \(\forall y \geqslant x\;\, \exists z \geqslant y: z \in V(p)\), then
  \(x \in V(p)\). This is what Humberstone
  (\citeproc{ref-Humberstone2011}{2011, 900}) calls
  \textbf{refinability}. It means that \(p\) only fails to be true at
  \(x\) if there is some refinement of \(x\) where it is settled as
  being untrue.
\end{itemize}

That's enough to provide a theory of truth for the language of
propositional logic, and Humberstone proposes the following clauses.

\begin{align*}
[\text{Vbls}] \quad & \mathcal{M} \models_x p_i \text{ iff } x \in V(p_i); \\
[\neg] \quad & \mathcal{M} \models_x \neg A \text{ iff } \forall y \geqslant x, \, \mathcal{M} \nmodels_y  A; \\
[\wedge] \quad & \mathcal{M} \models_x A \wedge B \text{ iff } \mathcal{M} \models_x A \text{ and } \mathcal{M} \models_x B; \\
[\vee] \quad & \mathcal{M} \models_x A \vee B \text{ iff } \forall y \geqslant x \;\, \exists z \geqslant y : \mathcal{M} \models_z A \text{ or } \mathcal{M} \models_z B. \\
[\rightarrow] \quad & \mathcal{M} \models_x A \rightarrow B \text{ iff } \mathcal{M} \models_x \neg(A \wedge \neg B)
\end{align*}

The first two clauses are familiar, and the third looks much like the
intuitionist treatment of negation. The final clause amounts to treating
\(\rightarrow\) as a defined connective. The surprising one is
\([\vee]\). It says that a disjunction is true iff at any refinement, at
least one of the disjuncts could become true. This is not the most
obvious understanding of \(\vee\), but it allows Humberstone to prove
three striking results.

First, every sentence in the language is persistent. If
\(\mathcal{M} \models_x A\) and \(x \leqslant y\), then
\(\mathcal{M} \models_y A\). For any sentence, truth is always preserved
when moving to a refinement.

Second, refinability holds for all sentences in the language. Once we've
added negation to the language, we can restate refinability as the
condition that, for any \(A,\) if \(\mathcal{M} \nmodels_x A\), there is
some refinement \(y\) of \(x\) such that
\(\mathcal{M} \models_y \neg A\). Intuitively, if a sentence is not
true, then at some refinement its negation is true.

Third, if we define validity on possibility frames, written as
\(\vDash_{pf}\), as truth preservation, then we get a very classical
theory. It's not just that, within this language, \(\vdash_{CL}\) and
\(\vDash_{pf}\) exactly overlap. It's also that, in any expanded
language, as long as persistence and refinability hold, the introduction
and elimination rules for \(\vdash_{CL}\) also hold for \(\vDash_{pf}\).

It's useful to work through one of these rules to show how the system
works. Because it will play a role in what follows, I'll show that
disjunction elimination holds for \(\vDash_{pf}\). One natural
formulation\footnote{There are other formulations which make weaker
  structural assumptions, but this one will do to show that we're
  preserving classical logic.} of disjunction elimination says that
\(\Gamma \cup \{A\} \vdash_{CL} C\),
\(\Delta \cup \{B\} \vdash_{CL} C\), and \(\Phi \vdash_{CL} A \vee B\)
together imply \(\Gamma \cup \Delta \cup \Phi \vdash_{CL} C\). Let's
show that the same principle holds for \(\vDash_{pf}\).

Assume that \(\Gamma \cup \{A\} \vDash_{pf} C\),
\(\Delta \cup \{B\} \vDash_{pf} C\), and \(\Phi \vDash_{pf} A \vee B\),
and assume for reductio that for some possibility \(x\) where all
members of \(\Gamma \cup \Delta \cup \Phi\) are true,
\(\mathcal{M} \nmodels_x C\). By refinability, there is some
\(y \geqslant x\) such that \(\mathcal{M} \models_y \neg C\), i.e., for
all \(z \geqslant y,\) \(\mathcal{M} \nmodels_z C\). By persistence, all
members of \(\Gamma \cup \Delta \cup \Phi\) are true at \(y\). Since all
of \(\Phi\) is true at \(y\), we have
\(\mathcal{M} \models_y A \vee B\). By \([\vee]\), there is some
\(z \geqslant y\) such that \(\mathcal{M} \models_z A\) or
\(\mathcal{M} \models_z B\). By persistence again, all members of
\(\Gamma \cup \Delta\) are true at \(z\). If
\(\mathcal{M} \models_z A\), then all of \(\Gamma \cup \{A\}\) are true
at \(z\), so \(\mathcal{M} \models_z C\). If
\(\mathcal{M} \models_z B\), then all of \(\Delta \cup \{B\}\) is true
at \(z\), so \(\mathcal{M} \models_z C\). Either way,
\(\mathcal{M} \models_z C\). But \(z \geqslant y\) and
\(\mathcal{M} \models_y \neg C\), so \(\mathcal{M} \nmodels_z C\).
Contradiction.

Humberstone's system looks non-classical. It includes points that make
true neither \(A\) nor \(\neg A\) for some \(A\). It includes a rule for
negation that looks more like intuitionistic negation than classical
negation. But looks are misleading; the system is still classical. All
the sequents, and all the admissible rules on sequents, carry over from
classical logic. We have a non-standard semantics for a very standard
logic.

Possibility frames have attracted some attention recently as an
interesting classical alternative to possible worlds semantics. We've
already mentioned Holliday (\citeproc{ref-Holliday2025}{2025}), and his
work will come up extensively in what follows. Matthew Harrison-Trainor
(\citeproc{ref-HarrisonTrainor2019}{2019}) has shown how to extend the
semantics to first-order logic. Brian Weatherson
(\citeproc{ref-WeathersonHoPF}{2026}) applies possibility frame
semantics to a puzzle about counterfactuals. And Fabrizio Cariani
(\citeproc{ref-CarianiUnpub}{2025}) uses possibility frames to model the
open future. This paper extends that list; I'm going to use them to
provide a better way of talking about fictional truth.

Since we need modal operators, we need to say more about \(\mathbb{R}\).
First we'll talk about modals. We start with a familiar truth condition
for \(\Box_i\):

\begin{align*}
[\Box_i] \quad & \mathcal{M} \models_x \Box_i A \text{ iff } \forall y \, (x R_i y \Rightarrow \mathcal{M} \models_y A)
\end{align*}

In words, \(\Box_i A\) is true at \(x\) iff \(A\) is true at every
\(R_i\)-accessible possibility. Humberstone treats \(\Diamond\) as a
defined connective, \(\Diamond_i A\) just means \(\neg \Box_i \neg A\),
and I'll do the same.

That alone isn't enough to guarantee that boxed sentences are persistent
and refinable. There are a few ways to constrain how \(\mathbb{R}\) and
\(\leqslant\) relate to get that guarantee; I'll use one suggested by
Holliday. He calls this constraint \textbf{R-Win}.

\begin{description}
\tightlist
\item[\textbf{R-Win}]
\(xR_iy \leftrightarrow \forall y' \geqslant y \exists x' \geqslant x \forall x'' \geqslant x' \exists y'' \geqslant y': x''R_iy''\)
\end{description}

It is easy to show that if each \(R_i\) satisfies \textbf{R-Win}, then
even after adding \(\Box_i\) and \(\Diamond_i\) to the language, all
sentences will satisfy persistence and refinability. It is a little
stronger than what's needed for these conditions, but the extra strength
turns out to be a helpful simplification.

One important consequence of \textbf{R-Win} is that each \(R_i\) is
carried downwards along refinement. That is, if \(x R_i y\) and
\(y \leqslant y'\), then \(x R_i y'\). More colloquially, if \emph{a}
can access \emph{b}, \emph{a} can access any refinement of \emph{b}.
This will be important when we get to epistemic modality.

Holliday (\citeproc{ref-Holliday2025}{2025, sec. 6.3}) notes that one
upside of adding \textbf{R-Win} is that it means we get back \emph{some}
familiar results from the traditional possible worlds semantics for
modal logic. A possibility frame validates
\(\Box_i \Phi \rightarrow \Phi\) iff \(R_i\) is reflexive; it validates
\(\Box_i \Phi \rightarrow \Box_i \Box_i \Phi\) iff \(R_i\) is
transitive; and it validates \(\Diamond_i \top\) iff \(R_i\) is serial.
What's more important here, and what Holliday also notes, is that we
don't get the same parallel for all axioms.

In possible world semantics, the frame condition for
\(\Box_i \Phi \vee \Box_i \neg \Phi\), or equivalently,
\(\Diamond_i \Phi \rightarrow \Box_i \Phi\), is that each point accesses
at most one point. That is,
\(\forall x, y, z ((xRy \wedge xRz) \rightarrow y = z)\). But in a
possibility frame semantics, Holliday (\citeproc{ref-Holliday2025}{2025,
229}) shows that what it takes to validate
\(\Diamond_i \Phi \rightarrow \Box_i \Phi\) is (13):

\begin{enumerate}
\def\labelenumi{(\arabic{enumi})}
\setcounter{enumi}{12}
\tightlist
\item
  \(\forall x, y (xR_iy \rightarrow \exists x^\prime \geqslant x \forall z (x^\prime \geq z \rightarrow \exists x^{\prime\prime} (x^{\prime\prime} \geq y \wedge x^{\prime\prime} \geqslant z)))\)
\end{enumerate}

In Figure~\ref{fig-simple-frame}, I've displayed a simple frame that
validates \(\Box_i \Phi \vee \Box_i \neg \Phi\) without being
functional. The solid lines represent \(R\)-accessibility; the dashed
lines represent refinements. From \(x_1\) all the \(y_i\) are
accessible. So neither \(\Box p\) nor \(\Box \neg p\) are true at
\(x_1\). But the disjunction \(\Box p \vee \Box \neg p\) is true, since
every refinement chain coming out of \(x_1\) eventually validates either
\(\Box p\) (on the chain going to \(x_2\)) or \(\Box \neg p\) (on the
chain going to \(x_3\)).\footnote{A similar construction shows that we
  can add a counterfactual that validates Conditional Excluded Middle
  without assuming that for each possibility \(x\) and possibly true
  proposition \(\Phi\), there is a unique nearest possibility where
  \(\Phi\) is true.}

\begin{figure}

\centering{

\pandocbounded{\includegraphics[keepaspectratio]{new-work_files/figure-pdf/fig-simple-frame-1.jpg}}

}

\caption{\label{fig-simple-frame}Possibility frame with refinement
(dashed, \(\geqslant\)) and \(R\)-accessibility (solid arrows)
relations}

\end{figure}%

\section{Dempster-Shafer Functions}\label{sec-sec2}

Orthodox opinion in philosophy is that if a thinker's uncertainty can be
measured numerically, so there is some function \(Bel\) from
propositions to numbers representing the agent's degree of belief in
each proposition, \(Bel\) must be a probability function. This view is
usually called \emph{probabilism}. Assume, for simplicity for now, that
there is a finite set \(\mathcal{F}\) of maximally specific relevant
possibilities, and each proposition the thinker has an attitude towards
is a subset of \(\mathcal{F}\). Saying that \(Bel\) is a probability
function is equivalent to saying there is some mass function \(m\) from
singletons of members of \(\mathcal{F}\) to non-negative reals such
that:

\begin{enumerate}
\def\labelenumi{\arabic{enumi}.}
\tightlist
\item
  \(\sum_{x \in \mathcal{F}} m(\{x\}) = 1\)
\item
  \(Bel(A) = \sum_{x \in A} m(\{x\})\)
\end{enumerate}

The most widely endorsed alternative to probabilism is that \(Bel\) only
needs to satisfy the slightly weaker constraints of
\emph{Dempster-Shafer belief functions}
(\citeproc{ref-Dempster1967}{Dempster 1967};
\citeproc{ref-Shafer1976}{Shafer 1976}). I'll call these \emph{belief
functions} from here on. Belief functions are such that there is a mass
function \(m\) from non-empty subsets of \(\mathcal{F}\) to non-negative
reals such that:

\begin{enumerate}
\def\labelenumi{\arabic{enumi}.}
\tightlist
\item
  \(\sum_{X \subseteq \mathcal{F}} m(X) = 1\)
\item
  \(Bel(A) = \sum_{X \subseteq A} m(X)\)
\end{enumerate}

If \(m\) only assigns positive mass to singletons, then \(Bel\) will be
a probability function, so all probability functions are belief
functions. But the converse is not true. The simplest case is where
there are just two possibilities: \(a, b\), and \(m\) assigns mass
\(\frac{1}{3}\) to each of \(\{a\}, \{b\}\), and \(\{a, b\}\). Then
\(Bel(\{a\}) = Bel(\{b\}) = \frac{1}{3}\), while \(Bel(\{a, b\}) = 1\).
We could state the rules on belief functions axiomatically, but this
more semantic presentation is both more intuitive, and connects more
smoothly to the application in section \ref{sec-sec4}, where the mass
function will be put to work.

What motivates the more restrictive orthodox view, that \(Bel\) has to
be probabilistic? In the twentieth century, the primary motivations were
largely pragmatic: Dutch Book arguments, and arguments from the
representation of preferences over gambles. There were two problems with
these arguments. One was that they were pragmatic, and hence the wrong
kind of thing to ground a purely epistemic constraint. The other, as
Patrick Maher (\citeproc{ref-Maher1997}{1997}) pointed out, was that
they were question-begging against the most important opponent, the
proponent of belief functions. David Christensen
(\citeproc{ref-Christensen1996}{1996}) solved the first problem by
developing `depragmatised' Dutch Book arguments, but this (a) made the
other problem worse, and (b) did not generalise to provide constraints
on updating degrees of belief.

So these days the key argument for orthodoxy is some kind of accuracy
based argument. The canonical such argument, due to James Joyce
(\citeproc{ref-Joyce1998}{1998}), says that if \(Bel\) is not a
probability function, then for any plausible measure of the accuracy of
one's degrees of belief, there is some alternative \(Bel^\prime\) which
is guaranteed to be more accurate.

The \emph{guarantee} here is broadly modal. It means that \(Bel^\prime\)
is more accurate than \(Bel\) in every possible world, where worlds are
understood to be sets of sentences which are complete (every sentence is
either true or false) and closed under classical entailment. These
assumptions are load-bearing. J. Robert G. Williams
(\citeproc{ref-Williams2016prob}{2016}) shows that if we replace
classical entailment with intuitionistic entailment then the constraint
is not that \(Bel\) is a probability function, but that it's a certain
kind of intuitionistic probability function. Here we want to look at
what happens if we keep the idea that we're using classical entailment,
but drop the assumption of completeness.

In particular, consider what happens if the thinker knows they are in
some possibility in a possibility frame, but does not know which one.
For example, consider someone who knows that they are in one of the
\(x_i\) in Figure~\ref{fig-simple-frame}. To connect it to the earlier
discussion, let the \(R\) relation be accessibility for \(\Box_{ER}\),
the operator meaning \emph{It's true in Eleanor Rigby that}, and let
\(p\) be that The Couple has grandchildren. Then \(\Box_{ER} p\) is true
at \(x_2\), \(\neg \Box_{ER} p\) is true at \(x_3\), as is (almost
equivalently) \(\Box_{ER} \neg p\), while none of these are true at
\(x_1\). (Though both \(\Box_{ER} p \vee \neg \Box_{ER} p\) and
\(\Box_{ER} p \vee \Box_{ER} \neg p\) are true at \(x_1\).)

Now imagine the thinker has
\(Bel(\Box_{ER} p) = Bel(\neg \Box_{ER} p) = \frac{1}{3}\). At any
possible \emph{world}, the function \(Bel^\prime\), which has
\(Bel^\prime(\Box_{ER} p) = Bel^\prime(\neg \Box_{ER} p) = \frac{1}{2}\),
is guaranteed to be more accurate. But at \(x_1\), \(Bel\) is more
accurate than \(Bel^\prime\). After all, both \(\Box_{ER} p\) and
\(\neg \Box_{ER} p\) fail to be true, so it's more accurate to have
lower credences in them. And \(Bel\) is lower than \(Bel^\prime\) on
each proposition.

I conclude that an assumption that's needed to make the
accuracy-dominance argument for probabilism work is that every question
the thinker is uncertain about is guaranteed to be factual. That is, it
only works if the thinker is certain they are not in possibilities like
\(x_1\) where there is no fact of the matter about whether
\(\Box_{ER} p\) is true. (Or that they are in \(y_1\), where there is no
fact of the matter about whether \(p\) is true; but usually we have more
evidence we are not in such sparse possibilities as \(y_1\).)

There are two very natural responses the proponent of accuracy dominance
arguments might make here. One is to simply accept this result, and say
that their argument was only meant to apply to questions known to be
factual. That's still a lot of questions, and if we have to qualify our
probabilism when thinking about the non-factual, well there's still a
lot of factual questions about which probabilism is correct. If that's
the next move, the obvious next question is just how wide the realm of
the factual is known to be, but I'll set that aside here. An alternative
response is to argue that this restriction to the factual is no
restriction at all, since all questions which can be sensibly asked are
known to have factual answers. In the next section, I'll argue that
questions about fiction show this isn't right.

\section{Fictional Reality}\label{sec-sec3}

Let's go back to (3), repeated here.

\begin{enumerate}
\def\labelenumi{(\arabic{enumi})}
\setcounter{enumi}{2}
\tightlist
\item
  In the fiction, The Husband's name is Vincent.
\end{enumerate}

We can wonder about whether this is true. Intuitively, it is more likely
that his name is Vincent than that it's Vaibhav. We could get even
clearer about just how likely it is to be Vincent by looking up
frequency of names for newly married men around where and when the story
is set. (That's not to say we can specify those details; but it seems
clear the story is set somewhere in the UK, probably Liverpool or
Greater London, not long before 1966.)

If all questions we can wonder about are factual, there seem to be two
live options. One, defended by Lewis (\citeproc{ref-Lewis1978b}{1978}),
is that we know (3) is false. In the introduction I surveyed some
arguments from Blumberg and Holguı́n
(\citeproc{ref-BlumbergHolguin2025}{2025}) against that view. Here's one
more argument I didn't mention there. If we know (3) is false, then
everything I said in the last paragraph is false. It's not true that his
name is more likely Vincent than Vaibhav, or that the story is probably
set in Liverpool or Greater London. That is, at least, seriously
revisionary.

Blumberg and Holguín conclude from this that there is an unknown fact of
the matter about whether (3) is true. As they note, this view is widely
held to be ``absurd''. That's for the excellent reason that it is, in
fact, absurd! How could Paul McCartney writing a few lines settle all
the questions we could ask about the characters, from what grade Father
McKenzie got in his last high school English exam, to how many eggs The
Wife ate in her lifetime? What could ground those facts? Those intuitive
arguments seem compelling to me, but rather than rest on them, I'll note
that their view has fairly revisionary consequences for how we think
about the relationship between fictions.

These days, the bulk of fictions produced are fan fictions (or
\emph{fanfics}); stories produced by fans recycling characters from
established works. Sometimes these don't bother with consistency with
the underlying source material, as if the writer had put Julius Caesar
into the Korean War. But a large body of fanfics are, or at least
attempt to be, `canon-compliant'. That is, they aim to write stories
which do not contradict any truths in the source material. If Blumberg
and Holguín are right, then no two fanfics which contradict each other
can be canon-compliant. Since the source settles all questions, at most
one of the disagreeing fanfics can agree with the source.

This doesn't just affect fanfics; the same issue comes up with sequels.
It is usually taken to be a constraint on sequels that they should agree
with what was true in the prequel, or if not explain the divergence. If
Blumberg and Holguín are correct, approximately no sequels satisfy this
constraint. If we take \emph{A Study in Scarlet} to be a self-contained
fiction, then it is incredibly unlikely that the protagonist of it will
one day solve a murder in Dartmoor, have a fight to the death above the
Reichenbach Falls, protect a woman whose stepfather tries to murder her
by passing a snake through a heating vent, and so on. On their view,
there is a fact at the end of \emph{A Study in Scarlet} about which (if
any) of these things Holmes will one day do, and it seems improbable he
will do all of them. So it's not just Watson's war wound which changes
from one thing to another over the Holmes stories, it's Holmes's entire
life trajectory.

Perhaps we could deny the assumption that \emph{A Study in Scarlet} is
self-contained; what's true in a book which has a sequel is determined
by what happens in later books. That too, however, would be revisionary.
We don't think that the later descriptions of where the wound is get to
be correct solely in virtue of being later.

It's implausible that all fictions are complete in the sense that
Blumberg and Holguín suggest, and assuming that they are requires making
massive changes to how we think about the relationship between fictions.
So we should try something else. My suggestion is that the logic of
\emph{In the fiction} is a modal logic where the underlying semantics is
a possibility frame. For the reasons Blumberg and Holguín give, it's
plausible that this logic is functional, i.e., it satisfies
\(\Diamond \Phi \rightarrow \Box \Phi\). But as long as we use
possibility frames, this does not require that there is a particular
world that the fiction picks out.

Blumberg and Holguín are aware of the option that one could deny that
the fiction is complete in this way, and they offer a series of
objections to it. The most pressing is that if the fiction is
incomplete, then (14) will be true, when \emph{p} is replaced by any
claim that lacks a truth value in the fiction.

\begin{enumerate}
\def\labelenumi{(\arabic{enumi})}
\setcounter{enumi}{13}
\tightlist
\item
  I know whether \(p\) is true, but I don't know whether \(p\).
\end{enumerate}

That would be very bad, but fortunately it isn't a consequence of the
view that the fiction is incomplete. For (14) to be true, it would not
merely have to be the case that some claims in the fiction lack truth
values. It would also have to be the case that it is known that \(p\) is
one such proposition about which the fiction is incomplete. This isn't
just an aspect of the objection based on (14); all of Blumberg and
Holguín's objections make this further assumption.

Somewhat surprisingly, if we represent incompleteness with possibility
frames, then it is never knowable that the fiction is incomplete with
respect to \(p\). So this further assumption of the objection isn't
correct, and the view defended here is not committed to implausible
claims like (14). To see why this is so, we need to return to a point I
flagged earlier about the nature of epistemic modality in possibility
frames.

Assume that \emph{S knows} is a normal modal operator, so there is some
relation \(R_k\) such that \emph{S knows that} \(\Phi\) is true at \(x\)
iff for all \(y\) such that \(xR_ky\), \(\Phi\) is true at \(y\). Since
knowledge is factive, this \(R_k\) will have to be reflexive.\footnote{This
  isn't quite as obvious a point as it is in the possible worlds
  framework. But as noted above, Holliday proves that if we assume
  \textbf{R-Win}, then all and only reflexive frames validate
  \(\Box \Phi \rightarrow \Phi\).} If, as is standard, we assume that
\(xRy\) and \(y^\prime \geqslant y\) imply that \(xRy^\prime\), it
follows that whenever \(x^\prime \geqslant x\), then \(xR_kx^\prime\).
By refinability, if neither \(\Phi\) nor \(\neg \Phi\) are true at
\(x\), then there are successor points to \(x\) where \(\Phi\) is true,
and where \(\neg \Phi\) is true. Crucially, this means it won't be the
case that \emph{It is not true that} \(\Phi\) holds at all points \(y\)
such that \(xR_ky\). So it won't be \emph{known} that \(\Phi\) is not
true.

A special case of this is where \(\Phi\) is \(\Box_{ER}p\), for some
claim \(p\) that is intuitively unsettled by the text of \emph{Eleanor
Rigby}. There might be no fact of the matter about whether, in
\emph{Eleanor Rigby}, \emph{p} is true; but we can never know that. So,
among other things, (14) is never true, defanging the objection.

This might seem like a rather Pyrrhic victory, however. There are
several reasons we might worry about a view that says the fiction is
incomplete, but we can never know where it is incomplete. Here are three
immediate worries. First, this move looks ad hoc. Second, it looks
implausible, given how sparse \emph{Eleanor Rigby} is. Third, we might
worry that once we go down this route, it will be impossible to say what
is different about the possibility frame view as compared to Blumberg
and Holguín's fictional reality view. Let's take these in turn.

It's certainly true that if we simply stipulated that the fictional
can't be known to be incomplete in a respect, that would be implausibly
ad hoc. But that isn't what happened. The lack of knowledge follows from
general principles relating \(\geqslant\) and \(\Box_{ER}\), which
Humberstone and Holliday defend on independent grounds. It's a
surprising consequence of those principles, but the principles weren't
engineered to get this result.

The implausibility worry is more pressing. I would have thought offhand
that it's obviously true, and hence known, that \emph{Eleanor Rigby} is
incomplete in several respects. If we think possibility frames are the
best way to represent incompleteness, and they say this incompleteness
is unknowable, then maybe that's some argument against this intuition,
but hardly a strong one. The best response, I feel, is that (14) seems
obviously bad, but if there was a knowable incompleteness it would be
true, so maybe we should live with all incompleteness being unknowable.

There are two different worries about expressibility, and I'll just
address one here.

One worry, which Cariani (\citeproc{ref-CarianiUnpub}{2025}) addresses
at length, is that we can't add to the language a way of saying anything
is incomplete without messing up the persistence and refinability
conditions. If you could say, truly, that \(p\) is incomplete, that
would have to be true on all refinements, which it plainly can't be
given refinability. As Cariani notes, there are various ways to finesse
this using indexicals. Roughly, you can have something in the language
which means that \(p\) is \emph{actually} incomplete. The details here
are messy, but ultimately manageable. But they don't, I think, help with
the deeper worry about inexpressibility.

That worry is that we want some way of saying we are pretty sure that
there is no fact of the matter about whether The Couple have
grandchildren. In the next section, I'll address this concern. The short
version is that while we can't know there is no fact of the matter about
\(\Box_{ER}p\), we can rationally have arbitrarily low credence in both
\(\Box_{ER}p\) and \(\neg \Box_{ER}p\). What's distinctive about realms
featuring incompleteness is that these pairs of low credences can be
rational.

\section{Uncertainty about the Non-Factual}\label{sec-sec4}

Think back to Figure~\ref{fig-simple-frame}, our model for uncertainty
about \emph{Eleanor Rigby}. There are six possibilities in the frame,
and four of them are complete. Since we know that Eleanor Rigby is a
fictional character, and \(y_1, y_2\) and \(y_3\) are possibilities
where she is real, we know that none of those possibilities represent
reality. But the other three, \(x_1\), \(x_2\) and \(x_3\), are all
candidates for being reality. From the outside, we might think that any
assignment of credences to those three possibilities which sum to 1 will
be coherent. This is a fairly strong assumption, and I'll come back to
it presently; first, let's think about what happens if it is true.

If we again let \(p\) be that The Couple has grandchildren, then
\(Cr(p)\) will be \(Cr(\{x_2\})\), and \(Cr(\neg p)\) will be
\(Cr(\{x_3\})\). Since \(Cr(\{x_1\})\) could be positive, it is possible
that \(Cr(p) + Cr(\neg p) < 1\). Indeed, by expressing credences in
\(p\) and \(\neg p\) which are both low, we can express confidence that
we are in \(x_1\), i.e., that there is no fact of the matter about
\(p\). In this case, we really can whistle something we can't say. Even
if we can't say directly that there is (probably) no fact of the matter
about \(p\), we can express confidence that there is no such fact
indirectly, via our credences in \(p\) and \(\neg p\).

This all assumed that credences in \(x_1\), \(x_2\) and \(x_3\) should
sum to 1, and that it is permissible for credence in \(x_1\) to be
positive. Both claims are contentious, and they push us in opposite
directions. The arguments of Section~\ref{sec-sec3} are, in effect,
arguments that \(x_1\) might be actual, and hence that it is reasonable
to give it positive credence. Indeed, those arguments suggest it should
be given a rather large credence. As noted earlier, the best argument
against non-probabilistic credences is that they are accuracy dominated.
If there is guaranteed to be a fact of the matter about whether \(p\),
then the credence which assigns \(\frac{1}{3}\) to each of \(p\) and
\(\neg p\) will, in both possibilities, be overall less accurate than
the credence which assigns \(\frac{1}{2}\) to each. But if it is
possible that neither \(p\) nor \(\neg p\) is true, then the former
credence will be more accurate than the latter in that possibility. So
the standard argument against non-probabilistic credences won't go
through.

Still, we might ask whether there is some other way to show that this
credence is accuracy dominated. And further, once we've left the realm
of probabilities, we might wonder why I've insisted that the `outside'
credences in \(x_1\), \(x_2\) and \(x_3\) sum to 1. The answer to both
challenges comes from a result in Predd et al.
(\citeproc{ref-Predd2009}{2009}), building on earlier work by Finetti
(\citeproc{ref-DeFinetti1974}{1974}) and Lindley
(\citeproc{ref-Lindley1982}{1982}).

Take a non-empty sample space \(\Omega\) and a finite list
\(E_1, \ldots, E_n\) of subsets of it, which we call events. A
\emph{forecast} is a vector \(f \in
[0,1]^n\), one number per event. It is \emph{coherent} just in case
there is a probability measure \(\mu\) over \(\Omega\) with
\(f_i = \mu(E_i)\) for each \(i\). We do not assume that these events
form an algebra, and we do not assume that \(\Omega\) is finite, though
we do assume that the list is finite.

De Finetti provided a geometrical characterisation of coherence. For
each \(\omega \in \Omega\), let \(v_\omega\) record the truth values of
the events at \(\omega\), assigning 1 to each \(E_i\) containing
\(\omega\) and 0 to the rest. These are the possibly perfectly accurate
forecasts. Let \(V\) be the set of these possibly perfectly accurate
forecasts. Then \(f\) is coherent if and only if it lies in the convex
hull of \(V\). That is, it is coherent iff it is a weighted average of
the possibly perfectly accurate forecasts, with the weights all being
non-negative and summing to one.

What's relevant to us is a further result in Predd et al.~Score a
forecast at \(\omega\) using any additive strictly proper scoring rule.
Then a forecast is strongly dominated if some other forecast scores
better at every point of \(\Omega\). It is weakly dominated if some
other forecast scores at least as well everywhere, and better somewhere.
A coherent forecast is not weakly dominated; an incoherent one is
strongly dominated, and indeed strongly dominated by something
coherent.\footnote{The last point is not, I think, philosophically
  significant. It's bad to have accuracy dominated credences; it doesn't
  matter whether the dominating credences are coherent. When a game
  theorist deletes a strategy because it is dominated, they don't do any
  further checks about whether the dominating strategy is coherent, or
  undominated, or otherwise playable. That seems like the right
  approach; being dominated is a decisive consideration against doing
  something, no matter what other properties the dominating strategy
  has.} The coherent forecasts just are the ones that are not accuracy
dominated.\footnote{It is essential that the scoring rule be
  \emph{strictly} proper. Without it, as Predd et al.
  (\citeproc{ref-Predd2009}{2009}) show, we can have incoherent
  credences that are only weakly dominated. On how far propriety and
  additivity can be relaxed see Pruss (\citeproc{ref-Pruss2024}{2024}),
  and on infinite agendas see Kelley (\citeproc{ref-Kelley2023}{2023}).}

Now imagine that instead of just caring about whether The Couple have
grandchildren, we are interested in some longer list \(E_1, \dots, E_n\)
of questions about \emph{Eleanor Rigby}. Imagine that instead of
Figure~\ref{fig-simple-frame}, we have a much larger frame, one that
includes on the right-hand side all the ways these questions might, or
might not, be resolved. (The `might not' here is because we want
possibilities like \(y_1\) that leave some questions open.) I'm still
working with the assumption that possibilities like \(x_1\), which leave
some questions open, are genuine possibilities; they might be actual.
I'll also assume that all the items on the agenda are stateable
propositions. That is, they all are sets of possibilities which satisfy
persistence and refinability.

For any item \(A\) on the agenda, it is true at a possibility \(x\) iff
it is true at all the refinements of \(x\); that follows from
persistence and refinability. So if \(x\) is actual, the perfectly
accurate credence assigns 1 to \(A\) when \(x\)'s refinements settle
\(A\), and 0 otherwise. Now the key step is that the coherent credences,
the ones that are not accuracy dominated, are all and only those in the
convex hull of these possibly perfectly accurate credences. So a
coherent credal state distributes non-negative weights \(w_x\) (which
sum to 1) over the possibilities and sets \(Bel(A)\) to the total weight
on possibilities where \(A\) is true.

Given these weights, we can construct a mass function \(m\). For each
possibility \(x\), let \(T_x\) be the conjunction of all the truths at
\(x\). Then for any proposition \(A\), let \(m(A)\) be the sum of
\(w_x\) across the possibilities \(x\) such that \(T_x\) is equivalent
to \(A\). Normally there will be at most one such possibility. But it is
coherent for a frame to contain two possibilities at which exactly the
same things are true, so we need to define \(m(A)\) as a sum to capture
that. Either way, since the weights are non-negative and sum to 1, and
the masses are sums of weights within cells of a partition, the masses
are also non-negative and sum to 1. That's what they need to be in order
to be a mass function. Given \(m\), we can define
\(Bel(A) = \sum_{X \vDash A} m(X)\), where \(X\) ranges over
propositions. That is, \(Bel\) will be a Dempster-Shafer function, and
\(m\) its mass function.\footnote{This way of constructing \(Bel\) via
  \(m\) mirrors the original construction of Dempster-Shafer functions
  in Dempster (\citeproc{ref-Dempster1967}{1967}).} This gives the mass
function of section \ref{sec-sec2} an interpretation it lacked when I
first introduced it. We started there with a simple belief function,
where \(m\) assigns mass 1/3 to each of \(\{a\}\), \(\{b\}\) and
\(\{a, b\}\). Once that mass is turned into a belief function, it
represents the credences of a thinker who is (roughly speaking) equally
confident in Figure~\ref{fig-simple-frame} that actuality is \(x_1\),
\(x_2\) or \(x_3\). The key thing is that as long as any mass is
assigned to \(\{a, b\}\), i.e., any credence is given to \(x_1\) being
actual, then the resulting credences will be non-probabilistic.

That said, given a class of possible refinements, not any old belief
function can be defined on them. There can be propositions definable on
a model which are not equivalent to \(T_x\) for any possibility \(x\) in
the model. For instance, in Figure~\ref{fig-simple-frame}, there is no
possibility \(x\) such that \(T_x\) is equivalent to the proposition
\(\{x_2, y_2\}\). So while these constructions show (a) that some
non-probabilistic credences are coherent, and (b) only credences that
correspond to Dempster-Shafer functions are coherent, given a particular
possibility frame, we can't derive the further result that any
Dempster-Shafer function definable over the propositions in it could be
immune to accuracy dominance.

Several objections deserve brief attention. First, someone might say:
``This doesn't actually save Dempster and Shafer themselves. Their work
wasn't framed in anti-realist terms.'' That's correct, and I don't claim
otherwise. This is \emph{new work} for Dempster-Shafer functions. These
functions are useful for anti-realists, but anti-realism is not what
originally motivated their development. The point is that anti-realists
have good reason to work with belief functions, not that this was always
the intended application.

Second, the approach seems to rest on an unusual account of disjunction.
How can it be that \(p\) is not true, and \(\neg p\) is not true, yet
their disjunction is true? This does seem strange from a classical
perspective. But I'd argue that disjunction \emph{is} strange, and we've
underinvested in understanding it. It's curious that the philosophical
literature contains roughly one hundred times as many papers and books
on conditionals as on disjunction. If disjunction is philosophically
important, as it must be if it's central to how possibilities are
structured, we should think more carefully about it.\footnote{Humberstone's
  encyclopedic \emph{The Connectives} has just as much on disjunction as
  on conditionals, and an important section of the disjunction chapter
  discusses possibilities. I'm following this chapter both in taking
  disjunction to be more complicated than it first looks, and to be
  closely connected to possibilities.} Saying that a philosophical view
is wrong because it fails given a simple Boolean account of disjunction
is no more plausible than saying that a view fails because it is
inconsistent with treating conditionals as material implication.

Third, one might protest that this is all rather specific to fiction,
and in particular to trivial questions about fiction like whether The
Couple has grandchildren. But it covers much more than that. For one
thing, there are fairly substantive questions about particular fictions,
ones which are central to how we understand that fiction, where we
should at least take seriously the possibility that there is no fact of
the matter.

But more importantly, if this framework is viable, it has implications
well beyond fiction. In any situation where we want to model
uncertainty, while leaving open the possibility that there is no fact of
the matter, the framework could have uses. Here are two cases that seem
especially promising.

What credences should an agent have in various propositions about the
external world when they have absolutely no evidence whatsoever? If we
think they should be certain that realism is correct for all external
world propositions, i.e., they should be certain that any proposition
statable about the external world is either true or has a true negation,
then a familiar problem arises. Let \(E\) be the actual evidence that a
normal person has, and \(H\) some proposition about the future they know
that is not part of their evidence.\footnote{Obviously if Williamson
  (\citeproc{ref-Williamson1998}{1998}) is right and all one's knowledge
  is evidence, this is impossible. It doesn't seem plausible to me that
  propositions about the future are part of my evidence, so I'm happy
  rejecting this. See Goldman (\citeproc{ref-Goldman2009}{2009}) for
  more reasons for being sceptical that evidence and knowledge are quite
  so tightly linked.} Since this person knows \(H\), they know (or at
least could know) \(\neg E \vee H\). But all their evidence tells
against this; it is conclusive evidence for the negation of one
disjunct. So they must have known it a priori. This is weird since it is
contingent. This structure has been used in many arguments in
epistemology in recent times, most notably as the core of the objection
Roger White (\citeproc{ref-White2006}{2006}) makes to the form of
dogmatism defended by James Pryor (\citeproc{ref-Pryor2000}{2000}).

The framework here offers a distinctive kind of response to this case.
Perhaps a priori, we should be unsure whether there is even a fact of
the matter about \(E\) and about \(H\). We should give much of our
credence to a thin possibility where none of \(E, H, \neg E, \neg H\)
are determined to be true, and all of them could be true at refinements
of the possibility. On learning \(E\), we don't just move credence
around between fully realised worlds, we move credence from this
incomplete possibility to the ones where \(E\) is true. (And, among
them, primarily to ones where \(H\) is true.) This lets us retain the
intuitively plausible claims that knowledge of \(H\) is grounded in
\(E\), but that we did not have a priori justification to believe
contingencies like \(\neg E \vee H\).

(David Jehle and Brian Weatherson (\citeproc{ref-JehleWeatherson}{2012})
aim to resolve this problem using intuitionistic probability. That
response requires taking on more non-classical commitments than the
response here. Martin Martin Smith
(\citeproc{ref-MartinSmith2017}{2017}) has argued that Jehle and
Weatherson's approach should be rejected because it requires making more
anti-realist assumptions than even the intuitionist should accept.)

Another application of the framework is to moral anti-realism and moral
uncertainty. Michael Michael Smith
(\citeproc{ref-MichaelSmith2002}{2002}) has argued that some
anti-realists can't capture moral uncertainty. This framework isn't
going to particularly help solve the problem Smith raises for the
committed anti-realist; on this framework, credence is credence of
truth, so if you're committed to moral claims not having truth-values,
you can't assign them any positive credence. But while the framework
can't help the committed anti-realist, it can still be useful. One might
worry that the reasons Smith offers will create problems even for
someone who isn't sure whether realism or anti-realism is correct, and
mean that this person can't sensibly have credences representing their
moral uncertainty. What the framework shows is that it is possible to
coherently have credences in moral propositions while leaving open that
some moral claims might be neither true nor false. So while the
framework has trivial cases involving The Couple as its cleanest
applications, its potential applications are much much wider.

Finally, someone will ask: ``If all this is true, and not all credences
have to be probabilistic, why is probabilism so unreasonably effective?
Why does it work so well in thinking about everything from markets to
weather forecasts to sporting events?'' The answer is that probabilism
is unreasonably effective \emph{given realism}. In domains where we can
safely assume that a complete truth exists, probabilism works. The
accuracy dominance argument is perfectly sound. This mirrors a point
Graham Priest (\citeproc{ref-Priest1991}{1991}) has made about logic:
classical logic is the correct logic for reasoning about situations that
are known to be consistent. I say much the same: probabilism is the
correct credence framework for situations that are known to be complete.
But when completeness is not guaranteed, then it is permissible, perhaps
mandatory, to use non-probabilistic belief functions. Once we're in an
environment where realism is up for grabs, the accuracy dominance
argument does not rule out non-probabilistic belief functions (provided
they are really Dempster-Shafer belief functions).

\section{What is Classical Logic?}\label{sec-sec5}

I've argued that (a) we should model uncertainty about fictions using
possibility frames rather than just structures built from possible
worlds, and (b) if we do that, the right constraints on credences (to
ensure that they are not accuracy dominated) are given by
Dempster-Shafer belief functions rather than probabilities. You might
worry this isn't a particularly novel response. Williams
(\citeproc{ref-Williams2012grad}{2012}) already showed that if the
underlying logic is non-classical, then the accuracy dominance argument
motivates different constraints to classical probabilism. Indeed, he
shows that if the underlying logic is supervaluational, the right
constraints are given by Dempster-Shafer functions. So is the argument
here just a special case of the familiar result that accuracy dominance
arguments have distinct consequences in non-classical settings?

The answer to this question turns on the terminological question
\emph{What is classical logic?} To say this question is terminological
is not to say it's uninteresting. We could treat it as a relatively
boring question, answered either by stipulation or a survey of language
users. But we need not do this. If we answer it by thinking about what
we take to be essential to classical logic, then it turns out to reveal
both some interesting distinctions, and some useful features of the
possibility frames approach.

Humberstone (\citeproc{ref-Humberstone2011}{2011}) distinguishes three
ways of thinking about logics. First, we could identify a logic with its
theorems, what Humberstone calls the FMLA approach. Second, we could
identify a logic with its valid single-conclusion sequents, what
Humberstone calls the SET-FMLA approach. Finally, we could identify a
logic with its multi-conclusion sequents, what Humberstone calls the
SET-SET approach. We get three more ways of thinking about logics by
starting with these three categories (theorems, single-conclusion
sequents, and multi-conclusion sequents) and asking what rules the logic
licenses for generating new outputs (theorems or sequents) from familiar
ones. The result is six ways of asking whether a logic might or might
not be classical, as represented in Table~\ref{tbl-classical-open}.

\begin{longtable}[]{@{}rcc@{}}
\caption{Six respects in which a logic might agree with classical
logic}\label{tbl-classical-open}\tabularnewline
\toprule\noalign{}
& Output & Process \\
\midrule\noalign{}
\endfirsthead
\toprule\noalign{}
& Output & Process \\
\midrule\noalign{}
\endhead
\bottomrule\noalign{}
\endlastfoot
\textbf{FMLA} & ? & ? \\
\textbf{SET-FMLA} & ? & ? \\
\textbf{SET-SET} & ? & ? \\
\end{longtable}

In that table, the left-hand column represents whether a logic has the
same results as classical logic, the same theorems or axioms, and the
right-hand column represents whether it permits the same rules to get
there. The `Output' column just cares about statements within a
particular language. For simplicity, I'll here just focus on the
language of propositional logic. The `Process' column, however,
considers how the logic does when the language is expanded.

For a famous example of this, consider the argument against
supervaluationism in Williamson (\citeproc{ref-Williamson1994}{1994}).
He notes that supervaluationism rejects several important rules of
classical logic, such as \(\vee\)-elimination. There are two important
things to note about this argument. First, it is within the SET-FMLA
framework. It can't be in FMLA, because nothing like
\(\vee\)-elimination is a rule there. (Typically, the only rule is modus
ponens; everything else in propositional logic is done with axioms.) And
it can't be SET-SET because there we wouldn't have to even think about
processes. After all, supervaluationism rejects the classically valid
sequent \(\emptyset \vDash p, \neg p\). What's distinctive about
supervaluationism is that within SET-FMLA, we keep the classical
outcomes, but reject some classical processes. The second point is that
we have to expand the language, Williamson adds a \(D\) operator, to see
the problem. This is necessary because within the original language, we
know that the processes produce all and only the correct outcomes.

We can summarise these results about supervaluationism using the kind of
table we've already introduced. In Table~\ref{tbl-sv} the ways in which
supervaluationism agrees and disagrees with classical logic are
summarised.

\begin{longtable}[]{@{}rcc@{}}
\caption{Agreements and disagreements between classical logic and
supervaluationism}\label{tbl-sv}\tabularnewline
\toprule\noalign{}
& Output & Process \\
\midrule\noalign{}
\endfirsthead
\toprule\noalign{}
& Output & Process \\
\midrule\noalign{}
\endhead
\bottomrule\noalign{}
\endlastfoot
\textbf{FMLA} & \(\tick\) & \(\tick\) \\
\textbf{SET-FMLA} & \(\tick\) & \(\cross\) \\
\textbf{SET-SET} & \(\cross\) & \(\cross\) \\
\end{longtable}

The \(\cross\) on the right of the second row is the core of
Williamson's objection, and it does seem like a potent objection to
describing supervaluationism as \emph{classical}.

There are ten possible distributions of ticks and crosses to that
table.\footnote{This count assumes that the processes are sound and
  complete with respect to the outputs, that \(A\) is a theorem in
  classical FMLA iff \(\emptyset \vDash A\) is valid, and
  \(\Gamma \vDash \Delta\) is valid iff for some finite subset
  \(\Upsilon\) of \(\Delta\), \(\Gamma \vDash A\), where \(A\) is the
  disjunction of the members of \(\Upsilon\).} Classical logic itself
has ticks in all six cells; a paradigmatic non-classical logic like
intuitionistic logic has crosses in all six cells. It's an interesting
question to ask how many of the other eight combinations are satisfied
by reasonably natural logics. In several cases, I suspect the results
will be far from natural, but there are some interesting cases.

For instance, consider the logic (call it \(L\)) that you get by
starting with intuitionistic (propositional) logic, and restricting
attention to just the fragment of the language where the only
connectives are \(\neg\) and \(\wedge\). Then introduce \(\vee\) and
\(\rightarrow\) by stipulation:
\(A \vee B =_{df} \neg(\neg A \wedge \neg B)\),
\(A \rightarrow B =_{df} \neg(A \wedge \neg B)\). The theorems of \(L\)
are all and only the classical theorems. But \(L\) disagrees with
classical logic on some sequents, e.g., classical logic endorses
\(\neg \neg A \vDash A\). And given some expansions of the language,
\(L\) violates modus ponens. So the relationship between \(L\) and
classical logic is shown by Table~\ref{tbl-l}.

\begin{longtable}[]{@{}rcc@{}}
\caption{Agreements and disagreements between classical logic and
\(L\)}\label{tbl-l}\tabularnewline
\toprule\noalign{}
& Output & Process \\
\midrule\noalign{}
\endfirsthead
\toprule\noalign{}
& Output & Process \\
\midrule\noalign{}
\endhead
\bottomrule\noalign{}
\endlastfoot
\textbf{FMLA} & \(\tick\) & \(\cross\) \\
\textbf{SET-FMLA} & \(\cross\) & \(\cross\) \\
\textbf{SET-SET} & \(\cross\) & \(\cross\) \\
\end{longtable}

The logic of truth-preservation on possibility frames, call it \(PF\),
is much more classical than this. As Humberstone
(\citeproc{ref-Humberstone1981a}{1981}) shows, no matter how we expand
the language, all of the sequent-to-sequent rules of classical logic
hold. Of course, it disagrees with classical logic once we allow
multiple conclusions; like supervaluationism, \(PF\) does not validate
\(\emptyset \vDash p, \neg p\). So the agreements and disagreements are
summarised by Table~\ref{tbl-pf}.

\begin{longtable}[]{@{}rcc@{}}
\caption{Agreements and disagreements between classical logic and
\(PF\)}\label{tbl-pf}\tabularnewline
\toprule\noalign{}
& Output & Process \\
\midrule\noalign{}
\endfirsthead
\toprule\noalign{}
& Output & Process \\
\midrule\noalign{}
\endhead
\bottomrule\noalign{}
\endlastfoot
\textbf{FMLA} & \(\tick\) & \(\tick\) \\
\textbf{SET-FMLA} & \(\tick\) & \(\tick\) \\
\textbf{SET-SET} & \(\cross\) & \(\cross\) \\
\end{longtable}

The key difference between supervaluationism and \(PF\) is on the
SET-FMLA row; \(PF\) does, and supervaluationism does not, respect all
classical processes. This seems significant. Think back to why
Williamson's results about \(\vee\)-elimination were dialectically
important. It was common knowledge that supervaluationism did not agree
with classical logic in SET-SET. It doesn't take much to see that it
doesn't validate \(\emptyset \vDash p, \neg p\). Arguably, that was part
of the point. Still, it was thought to be a way of preserving classical
logic without bivalence. The failure of \(\vee\)-elimination threatens
that. This line of reasoning only makes sense if what's really necessary
for a logic being classical is that it has ticks on the top four cells
of tables like Table~\ref{tbl-pf}. And that's exactly what \(PF\) does.

This is one reason for identifying classical \emph{logic} with its
outputs and processes in SET-FMLA. There are other familiar reasons. You
rarely see multiple conclusion logics in introductory, or even
intermediate, textbooks. You practically never see multiple conclusion
arguments used in philosophy outside logic. It would be a very strange
metaphysics paper that included an argument with P1, P2, P3, followed by
C1 and C2 meaning that at least one of C1 and C2 are correct.

A natural formulation of current practice is that classical \emph{logic}
is identified with something in SET-FMLA, while the standard, one could
say classical, \emph{semantics} for that logic is bivalent. So for a
logic to agree with both classical logic and semantics, it has to have
ticks in at least one (and probably both) of the cells on the bottom row
of Table~\ref{tbl-pf}. Of course, \(PF\) does not have that; it is not a
classical semantics. But because it has ticks in the four cells above
that, it is a classical logic. That is, it is a non-classical semantics
for a classical logic.

If we agree with this way of speaking, it is easy to say what is
distinctive about the possibility frames approach to uncertainty. It was
already known that we needed to assume classical logic to get an
accuracy dominance argument for probabilism. Indeed, Williams
(\citeproc{ref-Williams2012grad}{2012}) showed just how strong the
agreement with classical logic needed to be; the pattern we see in
Table~\ref{tbl-sv} wasn't enough. What's new here is that even a system
that has the same introduction and elimination rules as classical logic
is still not strong enough to ground the accuracy dominance argument.
That argument really needs classical semantics. That is, it needs
bivalence.

\section{Overview}\label{overview}

Fictional works, like \emph{Eleanor Rigby}, typically raise questions
that they don't answer. I've argued that some of these questions really
don't have an answer. It's not that it is indeterminate what the answer
is, or (as Blumberg and Holguín argue) that we are ignorant of the true
answer. There is no fact of the matter about how to fill in the story in
some respects. Using possibility frames lets us see how this can be the
case even if the logic of \emph{true in fiction F} is functional in the
sense that it validates \(\Diamond A \rightarrow \Box A\).

Once we are in a realm where there is no fact of the matter about some
questions, the accuracy-dominance argument for probabilism no longer
goes through. That argument presupposes that if one credence function is
more accurate than another given that \(A\) is true, and given that
\(\neg A\) is true, then it is more accurate full stop. Without a
background assumption of bivalence, there is a missing step here. And
when we are discussing uncertainty in a realm where there might not be a
fact of the matter about some questions, that assumption shouldn't be
granted.

But this doesn't mean anything goes for credences once we drop
bivalence. Assuming we are at some point in a possibility frame, then
the accuracy-dominance argument still gives us a substantive conclusion:
the credences that are not accuracy-dominated are Dempster-Shafer belief
functions. That's the new work I've suggested for these functions: they
are the right way to represent how a rational agent is uncertain about a
subject matter where they aren't sure the world is complete.

\protect\phantomsection\label{refs}
\begin{CSLReferences}{1}{0}
\bibitem[\citeproctext]{ref-BlumbergHolguin2025}
Blumberg, Kyle, and Ben Holguı́n. 2025. {``Fictional Reality.''}
\emph{Philosophical Review} 134 (2): 149--201. doi:
\href{https://doi.org/10.1215/00318108-11676103}{10.1215/00318108-11676103}.

\bibitem[\citeproctext]{ref-CarianiUnpub}
Cariani, Fabrizio. 2025. {``Modeling Future Indeterminacy in Possibility
Semantics.''} \url{https://philarchive.org/rec/CARMFI}.

\bibitem[\citeproctext]{ref-Christensen1996}
Christensen, David. 1996. {``Dutch-Book Arguments {D}e-Pragmatized:
Epistemic Consistency for Partial Believers.''} \emph{Journal of
Philosophy} 93 (9): 450--79. doi:
\href{https://doi.org/10.2307/2940893}{10.2307/2940893}.

\bibitem[\citeproctext]{ref-Dempster1967}
Dempster, Arthur. 1967. {``Upper and Lower Probabilities Induced by a
Multi-Valued Mapping.''} \emph{Annals of Mathematical Statistics} 38:
325--39. doi:
\href{https://doi.org/10.1214/aoms/1177698950}{10.1214/aoms/1177698950}.

\bibitem[\citeproctext]{ref-DeFinetti1974}
Finetti, Bruno de. 1974. \emph{Theory of Probability}. New York: Wiley.

\bibitem[\citeproctext]{ref-Goldman2009}
Goldman, Alvin. 2009. {``Williamson on Knowledge and Evidence.''} In
\emph{{Williamson on Knowledge}}, 73--91. Oxford: Oxford University
Press.

\bibitem[\citeproctext]{ref-HarrisonTrainor2019}
Harrison-Trainor, Matthew. 2019. {``First-Order Possibility Models and
Finitary Completeness Proofs.''} \emph{Review of Symbolic Logic} 12 (4):
637--62. doi:
\href{https://doi.org/10.1017/S1755020319000418}{10.1017/S1755020319000418}.

\bibitem[\citeproctext]{ref-Holliday2025}
Holliday, Wesley H. 2025. {``Possibility Frames and Forcing for Modal
Logic.''} \emph{Australasian Journal of Logic} 22 (2): 44--288. doi:
\href{https://doi.org/10.26686/ajl.v22i2.5680}{10.26686/ajl.v22i2.5680}.

\bibitem[\citeproctext]{ref-Humberstone1981a}
Humberstone, Lloyd. 1981. {``From Worlds to Possibilities.''}
\emph{Journal of Philosophical Logic} 10 (3): 313--39. doi:
\href{https://doi.org/10.1007/BF00293423}{10.1007/BF00293423}.

\bibitem[\citeproctext]{ref-Humberstone2011}
---------. 2011. \emph{The Connectives}. Cambridge, MA: MIT Press.

\bibitem[\citeproctext]{ref-JehleWeatherson}
Jehle, David, and Brian Weatherson. 2012. {``Dogmatism, Probability and
Logical Uncertainty.''} In \emph{New Waves in Philosophical Logic},
edited by Greg Restall and Gillian Russell, 95--111. Bassingstoke:
Palgrave Macmillan.

\bibitem[\citeproctext]{ref-Joyce1998}
Joyce, James M. 1998. {``A Non-Pragmatic Vindication of Probabilism.''}
\emph{Philosophy of Science} 65 (4): 575--603. doi:
\href{https://doi.org/10.1086/392661}{10.1086/392661}.

\bibitem[\citeproctext]{ref-Kelley2023}
Kelley, Mikayla. 2023. {``On Accuracy and Coherence with Infinite
Opinion Sets.''} \emph{Philosophy of Science} 90 (1): 92--128. doi:
\href{https://doi.org/10.1017/psa.2021.37}{10.1017/psa.2021.37}.
Preprint at arXiv:2007.14490.

\bibitem[\citeproctext]{ref-Lewis1978b}
Lewis, David. 1978. {``Truth in Fiction.''} \emph{American Philosophical
Quarterly} 15 (1): 37--46. Reprinted in his \emph{Philosophical Papers},
Volume 1, Oxford: Oxford University Press, 1983, 261-275. References to
reprint.

\bibitem[\citeproctext]{ref-Lindley1982}
Lindley, D. V. 1982. {``Scoring Rules and the Inevitability of
Probability.''} \emph{International Statistical Review} 50: 1--26.

\bibitem[\citeproctext]{ref-Maher1997}
Maher, Patrick. 1997. {``Depragmatised Dutch Book Arguments.''}
\emph{Philosophy of Science} 64 (2): 291--305. doi:
\href{https://doi.org/10.1086/392552}{10.1086/392552}.

\bibitem[\citeproctext]{ref-Predd2009}
Predd, Joel B., Robert Seiringer, Elliott H. Lieb, Daniel N. Osherson,
H. Vincent Poor, and Sanjeev R. Kulkarni. 2009. {``Probabilistic
Coherence and Proper Scoring Rules.''} \emph{IEEE Transactions on
Information Theory} 55 (10): 4786--92. doi:
\href{https://doi.org/10.1109/TIT.2009.2027573}{10.1109/TIT.2009.2027573}.
Preprint at arXiv:0710.3183.

\bibitem[\citeproctext]{ref-Priest1991}
Priest, Graham. 1991. {``Minimally Inconsistent {LP}.''} \emph{Studia
Logica} 50 (2): 321--31. doi:
\href{https://doi.org/10.1007/BF00370190}{10.1007/BF00370190}.

\bibitem[\citeproctext]{ref-Pruss2024}
Pruss, Alexander R. 2024. {``Necessary and Sufficient Conditions for
Domination Results for Proper Scoring Rules.''} \emph{Review of Symbolic
Logic} 17 (1): 132--43. doi:
\href{https://doi.org/10.1017/S1755020323000035}{10.1017/S1755020323000035}.

\bibitem[\citeproctext]{ref-Pryor2000}
Pryor, James. 2000. {``The Sceptic and the Dogmatist.''} \emph{No{û}s}
34 (4): 517--49. doi:
\href{https://doi.org/10.1111/0029-4624.00277}{10.1111/0029-4624.00277}.

\bibitem[\citeproctext]{ref-Shafer1976}
Shafer, Glenn. 1976. \emph{A Mathematical Theory of Evidence}.
Princeton: Princeton University Press.

\bibitem[\citeproctext]{ref-MartinSmith2017}
Smith, Martin. 2017. {``Intuitionistic Probability and the {B}ayesian
Objection to Dogmatism.''} \emph{Synthese} 194 (10): 3997--4009. doi:
\href{https://doi.org/10.1007/s11229-016-1120-2}{10.1007/s11229-016-1120-2}.

\bibitem[\citeproctext]{ref-MichaelSmith2002}
Smith, Michael. 2002. {``Evaluation, Uncertainty and Motivation.''}
\emph{Ethical Theory and Moral Practice} 5 (3): 305--20. doi:
\href{https://doi.org/10.1023/A:1019675327207}{10.1023/A:1019675327207}.

\bibitem[\citeproctext]{ref-Standefer2018}
Standefer, Shawn. 2018. {``Proof Theory for Functional Modal Logic.''}
\emph{Studia Logica} 106 (1): 49--84. doi:
\href{https://doi.org/10.1007/s11225-017-9725-0}{10.1007/s11225-017-9725-0}.

\bibitem[\citeproctext]{ref-WeathersonHoPF}
Weatherson, Brian. 2026. {``Humberstone on Possibility Frames.''}
\url{https://brian.weatherson.org/quarto/posts/possibility-frames/HoPF.html}.

\bibitem[\citeproctext]{ref-White2006}
White, Roger. 2006. {``Problems for Dogmatism.''} \emph{Philosophical
Studies} 131 (3): 525--57. doi:
\href{https://doi.org/10.1007/s11098-004-7487-9}{10.1007/s11098-004-7487-9}.

\bibitem[\citeproctext]{ref-Williams2012grad}
Williams, J. Robert G. 2012. {``Gradational Accuracy and Nonclassical
Semantics.''} \emph{Review of Symbolic Logic} 5 (4): 513--37. doi:
\href{https://doi.org/10.1017/S1755020312000214}{10.1017/S1755020312000214}.

\bibitem[\citeproctext]{ref-Williams2016prob}
---------. 2016. {``Probability and Nonclassical Logic.''} In \emph{The
Oxford Handbook of Probability and Philosophy}, edited by Alan Hájek and
Christopher Hitchcock, 248--76. Oxford University Press. doi:
\href{https://doi.org/10.1093/oxfordhb/9780199607617.013.12}{10.1093/oxfordhb/9780199607617.013.12}.

\bibitem[\citeproctext]{ref-Williamson1994}
Williamson, Timothy. 1994. \emph{{Vagueness}}. New York: Routledge.

\bibitem[\citeproctext]{ref-Williamson1998}
---------. 1998. {``Conditionalizing on Knowledge.''} \emph{The British
Journal for the Philosophy of Science} 49 (1): 89--121. doi:
\href{https://doi.org/10.1093/bjps/49.1.89}{10.1093/bjps/49.1.89}.

\end{CSLReferences}



\noindent Draft for discussion.


\end{document}
